% !TeX root = ../../Book4.tex
% 纲要标题：### 3.3 伴随函子
% 纲要位置：计划纲要.md，第 2519 行。

\section{伴随函子}
\label{b4:sec:0303}

% 正文按纲要写入；各节的基础结果在本章证明。

% 纲要内容（仅作写作计划，尚非教材正文）：
%
% * Hom 集伴随及单位、余单位
% * 三角恒等式
% * 自由—忘却及 Hom–张量伴随
%

\subsection{Hom 集伴随及单位、余单位}
\label{b4:subsec:0303:01}

自由模的可表性针对一个固定集合 \(X\)。当 \(X\) 也变化时，表示对象 \(R^{(X)}\) 本身组成函子，泛性质就同时联系两个范畴。伴随正是这种双变量的自然对应；它不要求两个范畴等价，也不要求某一方忘掉的全部信息都能恢复。

\begin{definition}\label{b4:def:adjunction}
设 \(\mathcal C,\mathcal D\) 为局部小范畴，\(F:\mathcal C\rightarrow\mathcal D\)、\(G:\mathcal D\rightarrow\mathcal C\) 为函子。若存在关于 \(X\in\mathcal C\)、\(Y\in\mathcal D\) 自然的双射
\[
\Phi_{X,Y}:\operatorname{Hom}_{\mathcal D}(FX,Y)
\xrightarrow{\sim}\operatorname{Hom}_{\mathcal C}(X,GY),
\]
则称 \((F,G,\Phi)\) 为一个\term[bansui]{伴随}[adjunction]，称 \(F\) 为 \(G\) 的\term[zuobansui]{左伴随}[left adjoint]，\(G\) 为 \(F\) 的\term[youbansui]{右伴随}[right adjoint]，记 \(F\dashv G\)。这里两变量自然性明确要求：对 \(a:X'\rightarrow X\)、\(b:Y\rightarrow Y'\) 及 \(h:FX\rightarrow Y\)，有
\begin{equation}\label{b4:eq:adjunction-naturality}
\Phi_{X',Y'}\bigl(b\circ h\circ F(a)\bigr)
=G(b)\circ\Phi_{X,Y}(h)\circ a.
\end{equation}
\end{definition}
\symidx{\(F\dashv G\)：\(F\) 为 \(G\) 的左伴随}

固定 \(X\) 后，\(FX\) 表示函子 \(Y\mapsto\operatorname{Hom}_{\mathcal C}(X,GY)\)；固定 \(Y\) 后，\(GY\) 表示反变函子 \(X\mapsto\operatorname{Hom}_{\mathcal D}(FX,Y)\)。一个伴随把这些表示按另一个变量自然地组织起来，不能只对每对对象随意选取集合双射。

\ProseParagraphBreak
两个恒等态射给出特别重要的映射：
\begin{equation}\label{b4:eq:adjunction-unit-counit-def}
\eta_X\coloneqq\Phi_{X,FX}(1_{FX}):X\longrightarrow GFX,
\qquad
\varepsilon_Y\coloneqq\Phi^{-1}_{GY,Y}(1_{GY}):FGY\longrightarrow Y.
\end{equation}
它们分别称为伴随的\term[bansui-danwei]{单位}[unit]与\term[bansui-yudanwei]{余单位}[counit]。单位将一个对象送入“先作 \(F\) 再作 \(G\)”得到的对象；余单位则从“先作 \(G\) 再作 \(F\)”的对象回到原对象。

\begin{proposition}\label{b4:prop:adjunction-transpose-formulas}
设 \(F:\mathcal C\rightarrow\mathcal D\)、\(G:\mathcal D\rightarrow\mathcal C\) 是局部小范畴之间的伴随 \(F\dashv G\)，其自然双射为 \(\Phi\)。由\eqref{b4:eq:adjunction-unit-counit-def}定义的映射组成自然变换 \(\eta:1_{\mathcal C}\Rightarrow GF\) 与 \(\varepsilon:FG\Rightarrow1_{\mathcal D}\)，且
\begin{equation}\label{b4:eq:adjunction-transpose-formulas}
\Phi_{X,Y}(h)=G(h)\circ\eta_X,
\qquad
\Phi^{-1}_{X,Y}(k)=\varepsilon_Y\circ F(k)
\end{equation}
对所有 \(h:FX\rightarrow Y\)、\(k:X\rightarrow GY\) 成立。
\end{proposition}
\begin{proof}
对 \(a:X\rightarrow X'\)，把 \(F(a):FX\rightarrow FX'\) 分别看作 \(F(a)\circ1_{FX}\) 与 \(1_{FX'}\circ F(a)\)，双射的自然性给出
\[
GF(a)\circ\eta_X
=\Phi_{X,FX'}\bigl(F(a)\bigr)
=\eta_{X'}\circ a.
\]
所以 \(\eta\) 自然。由于自然同构的逆仍自然，对 \(b:Y\rightarrow Y'\)，同样在 \(\Phi^{-1}\) 下比较 \(G(b)=G(b)\circ1_{GY}=1_{GY'}\circ G(b)\)，得到
\[
b\circ\varepsilon_Y=\varepsilon_{Y'}\circ FG(b).
\]
这证明 \(\varepsilon\) 自然。最后，把 \(h\) 写成 \(h\circ1_{FX}\)，在第二变量使用自然性，得到第一条转置公式；把 \(k\) 写成 \(1_{GY}\circ k\)，对逆双射在第一变量使用自然性，得到第二条公式。
\end{proof}

因此，每个 \(k:X\rightarrow GY\) 都唯一分解为 \(G(h)\eta_X\)，其中 \(h:FX\rightarrow Y\)；每个 \(h:FX\rightarrow Y\) 都唯一写成 \(\varepsilon_YF(k)\)。这分别是单位与余单位的泛性质。第二卷第6.3节中“指定纯张量的像”和“固定一个变量”，正是这两个公式在模范畴中的表现。

\subsection{三角恒等式}
\label{b4:subsec:0303:02}

从单位和余单位能否恢复整个伴随？两条转置公式已经给出候选映射；所需条件就是保证它们互逆。称下列两式为\term[sanjiao-hengdengshi]{三角恒等式}[triangle identities]：
\begin{equation}\label{b4:eq:triangle-identities}
\varepsilon_{FX}\circ F(\eta_X)=1_{FX},
\qquad
G(\varepsilon_Y)\circ\eta_{GY}=1_{GY}.
\end{equation}
每条都是从一个对象出发，经三重函子复合再回到原对象的复合等式。

\begin{theorem}[伴随的单位与余单位刻画]\label{b4:thm:adjunction-unit-counit}
设 \(\mathcal C,\mathcal D\) 为局部小范畴，\(F:\mathcal C\rightarrow\mathcal D\)、\(G:\mathcal D\rightarrow\mathcal C\) 为函子。关于两变量自然的伴随双射
\[
\operatorname{Hom}_{\mathcal D}(FX,Y)
\cong\operatorname{Hom}_{\mathcal C}(X,GY)
\]
与满足\eqref{b4:eq:triangle-identities}的自然变换对
\(\eta:1_{\mathcal C}\Rightarrow GF\)、\(\varepsilon:FG\Rightarrow1_{\mathcal D}\)
一一对应。由双射得到变换的公式是\eqref{b4:eq:adjunction-unit-counit-def}，反向得到双射及其逆的公式是\eqref{b4:eq:adjunction-transpose-formulas}。
\end{theorem}
\begin{proof}
先给定伴随双射。单位、余单位的自然性已经由\cref{b4:prop:adjunction-transpose-formulas}证明。将第一条转置公式所得的 \(\eta_X=\Phi(1_{FX})\) 再代入逆公式，得到
\[
1_{FX}=\Phi^{-1}(\eta_X)=\varepsilon_{FX}F(\eta_X).
\]
将 \(\varepsilon_Y=\Phi^{-1}(1_{GY})\) 代回正向公式，得到
\(1_{GY}=G(\varepsilon_Y)\eta_{GY}\)。所以三角恒等式成立。

反过来，设 \(\eta,\varepsilon\) 自然且满足三角恒等式，定义
\[
\Phi(h)=G(h)\eta_X,\qquad \Psi(k)=\varepsilon_YF(k).
\]
对于 \(k:X\rightarrow GY\)，由 \(\eta\) 关于 \(k\) 的自然性以及第二条三角恒等式，有
\[
\Phi\Psi(k)
=G(\varepsilon_Y)GF(k)\eta_X
=G(\varepsilon_Y)\eta_{GY}k=k.
\]
对于 \(h:FX\rightarrow Y\)，由 \(\varepsilon\) 关于 \(h\) 的自然性以及第一条三角恒等式，有
\[
\Psi\Phi(h)
=\varepsilon_YFG(h)F(\eta_X)
=h\varepsilon_{FX}F(\eta_X)=h.
\]
因此 \(\Phi,\Psi\) 互逆。再对 \(a:X'\rightarrow X\)、\(b:Y\rightarrow Y'\) 计算
\[
\begin{aligned}
\Phi\bigl(bhF(a)\bigr)
&=G(b)G(h)GF(a)\eta_{X'}\\
&=G(b)G(h)\eta_Xa
=G(b)\Phi(h)a,
\end{aligned}
\]
证明双射在两变量上自然。由新双射在恒等态射上的取值，重新得到的单位与余单位仍为 \(\eta,\varepsilon\)。另一方向则由已证转置公式恢复原双射，故两种构造互逆。
\end{proof}

可对照 Weibel 的伴随定义及单位、余单位定理\cite[A.6.1; Theorem A.6.2 and Exercise A.6.2]{Weibel1994HomologicalAlgebra}。上面的两个方向在正文均已证明；其中自然性用于移动 \(\eta\) 与 \(\varepsilon\)，三角恒等式再消去中间的复合，两者各有职责。

\ProseParagraphBreak
三角恒等式不意味着 \(\eta_X\) 与 \(\varepsilon_Y\) 互逆：它们通常连定义域、陪域都不同。例如自由交换群伴随中，一个集合被送到它的形式整数线性组合；回到具体群时，余单位把这些形式和求值。这个过程通常带来许多关系，远非对象不变的同构。

\subsection{自由—忘却及 Hom–张量伴随}
\label{b4:subsec:0303:03}

设 \(R\) 为含幺环。自由左模函子 \(F_R:\mathbf{Set}\rightarrow R\text{-}\mathbf{Mod}\) 将 \(X\) 送到 \(R^{(X)}\)，并把函数 \(a:X\rightarrow X'\) 送到唯一满足 \(e_x\mapsto e_{a(x)}\) 的模同态。恒等与复合律由各基向量的像确定。\secref{b4:sec:0301}给出的双射不仅在模变量上自然，在集合变量上也自然：对 \(f:R^{(X')}\rightarrow M\)，前复合 \(F_R(a)\) 对应函数 \(x\mapsto f(e_{a(x)})\)。所以
\[
F_R\dashv U.
\]
单位为 \(\eta_X(x)=e_x\)，余单位为
\[
\varepsilon_M:R^{(U(M))}\longrightarrow M,
\qquad\sum_mr_me_m\longmapsto\sum_mr_mm.
\]
这里 \(e_m\) 是形式符号；它不是直接把 \(m\) 当作自由模中的元素。第一条三角恒等式把 \(e_x\) 先送到 \(e_{e_x}\)，再送回 \(e_x\)；第二条将 \(m\) 先送到 \(e_m\)，再求值为 \(m\)。因此两式也可直接在生成元上看出。

\ProseParagraphBreak
另一组伴随来自固定双模。它是第二卷第6.3节 Hom–张量伴随定理的原有左模版本；本卷把它放入已建立的一般定义，不重新构造张量积。

\begin{proposition}\label{b4:prop:hom-tensor-adjunction}
设 \(R,S\) 为含幺环，\({}_SP_R\) 为双模。函子
\[
F=P\otimes_R-:R\text{-}\mathbf{Mod}\longrightarrow S\text{-}\mathbf{Mod},
\qquad
G=\operatorname{Hom}_S(P,-):S\text{-}\mathbf{Mod}\longrightarrow R\text{-}\mathbf{Mod}
\]
满足 \(F\dashv G\)，其中 \(G(N)\) 的左 \(R\) 作用为 \((rf)(p)=f(pr)\)。其单位与余单位为
\[
\eta_M(m)(p)=p\otimes m,
\qquad
\varepsilon_N(p\otimes f)=f(p).
\]
\end{proposition}
\begin{proof}
第二卷第6.3节的定理对任意左 \(R\) 模 \(M\) 和左 \(S\) 模 \(N\) 给出自然双射
\[
h:P\otimes_RM\rightarrow N
\quad\longleftrightarrow\quad
\bigl(m\mapsto(p\mapsto h(p\otimes m))\bigr).
\]
那里已用平衡关系证明良定义、左线性及两变量自然性；所用环和双模假设与本命题完全相同，故满足\cref{b4:def:adjunction}。将两个恒等映射代入便得到所列公式。这里再明确作用次序：
\[
\eta_M(rm)(p)=p\otimes rm=pr\otimes m
=\bigl(r\eta_M(m)\bigr)(p),
\]
所以单位左 \(R\) 线性；余单位的平衡性为 \(f(pr)=(rf)(p)\)，其左 \(S\) 线性为 \(f(sp)=sf(p)\)。第一条三角恒等式在纯张量上将 \(p\otimes m\) 送到 \(\eta_M(m)(p)=p\otimes m\)；第二条将 \(f\) 送到函数 \(p\mapsto\varepsilon_N(p\otimes f)=f(p)\)。因此这两个具体映射也直接满足一般定理的条件。
\end{proof}

在非交换情形中，\(N\) 不一定有左 \(R\) 作用，不能把 \((rf)(p)\) 改写为 \(rf(p)\)。\(P\) 的右 \(R\) 作用经输入端传到 Hom 上，正是该伴随能够成立的原因。

\begin{example}[交换化与包含函子]\label{b4:exm:abelianization-adjunction}
对群 \(H\)，记 \(H^{\mathrm{ab}}=H/[H,H]\)，其中 \([H,H]\) 是由交换子生成的正规子群。每个到交换群 \(A\) 的群同态都消去交换子，因而唯一通过商群分解，给出
\[
\operatorname{Hom}_{\mathbf{Ab}}(H^{\mathrm{ab}},A)
\cong\operatorname{Hom}_{\mathbf{Grp}}(H,A).
\]
前复合群同态或后复合交换群同态都与商映射分解相容，故交换化是包含函子 \(\mathbf{Ab}\rightarrow\mathbf{Grp}\) 的左伴随。单位为商映射 \(H\rightarrow H^{\mathrm{ab}}\)，余单位为 \(A^{\mathrm{ab}}\xrightarrow{\sim}A\)。非交换群的单位不是单射，而余单位总是同构。这也说明伴随不要求单位与余单位具有相同的可逆性。
\end{example}
